% songs/raubritter.tex
\begin{song}{Raubritter}{Volker und Georg Hausmann}{Capo -}

\cl{  Am              F         Am    E     Am}
\slbl{1.}Von der Festung gröhlt's, derbe Männerstimmen.
\cl{       F           Am       E      Am}
Rauher Kehle Sang, hell die Becher klingen.

\slbl{Refrain:}
\cl{G   C            G       Am               E}
Beherrschen dies Gebiet, singen stolz ihr Lied.
\cl{Am                               G   Am  G Am}
Raubritter, Raubritter, wie weit ist unser Land
\cl{F                                 C   F   G  Am}
Raubritter, Raubritter, wie stark ist uns're Hand

\cl{  Am                      F    Am      E   Am}
\slbl{2.}In uns'rer Knechtschaft Zeit griffen wir zu Waffen,
\cl{            F      Am            E        Am}
verdroschen uns're Herrn, Grafen und auch Pfaffen.

\slbl{Refrain:}
\cl{G   C            G       Am               E}
Beherrschen dies Gebiet, singen stolz ihr Lied.
\cl{Am                              G   Am  G Am}
Raubritter, Raubritter, so weit ist unser Land
\cl{F                                C   F   G  Am}
Raubritter, Raubritter, so stark ist uns're Hand

\cl{  Am              F      Am      E   Am}
\slbl{3.}Groß ist uns're Macht, solange wir vereint.
\cl{          F      Am            E     Am}
Nie fällt uns're Burg, trutzet jedem Feind.

\slbl{Refrain:}
\cl{G   C            G       Am               E}
Beherrschen dies Gebiet, singen stolz ihr Lied.
\cl{Am                              G   Am  G Am}
Raubritter, Raubritter, so weit ist unser Land
\cl{F                                C   F   G  Am}
Raubritter, Raubritter, so stark ist uns're Hand
\anmerkung{1967 geschrieben von Volker (Anti) und Georg (Nathan) Hausmann, wurde es das Ordenslied des „Orden der Raubritter", der sich vom Orden der Rabenklaue im Nerother Wandervogel abgespalten hat. \\Raubritter waren Mitglieder des Ritterstandes, die sich durch Fehden und Plünderungszüge bereicherten, obwohl Könige und Fürsten dies verboten. Raubritterburgen standen deswegen oft in Feindschaft mit großen Städten im Umkreis.}
\end{song}
